% Options for packages loaded elsewhere
\PassOptionsToPackage{unicode}{hyperref}
\PassOptionsToPackage{hyphens}{url}
\PassOptionsToPackage{dvipsnames,svgnames,x11names}{xcolor}
\documentclass[
  10pt,
  fontset=fandol]{ctexart}
\usepackage{xcolor}
\usepackage[margin=16mm]{geometry}
\usepackage{amsmath,amssymb}
\setcounter{secnumdepth}{-\maxdimen} % remove section numbering
\usepackage{iftex}
\ifPDFTeX
  \usepackage[T1]{fontenc}
  \usepackage[utf8]{inputenc}
  \usepackage{textcomp} % provide euro and other symbols
\else % if luatex or xetex
  \usepackage{unicode-math} % this also loads fontspec
  \defaultfontfeatures{Scale=MatchLowercase}
  \defaultfontfeatures[\rmfamily]{Ligatures=TeX,Scale=1}
\fi
\usepackage{lmodern}
\ifPDFTeX\else
  % xetex/luatex font selection
\fi
% Use upquote if available, for straight quotes in verbatim environments
\IfFileExists{upquote.sty}{\usepackage{upquote}}{}
\IfFileExists{microtype.sty}{% use microtype if available
  \usepackage[]{microtype}
  \UseMicrotypeSet[protrusion]{basicmath} % disable protrusion for tt fonts
}{}
\makeatletter
\@ifundefined{KOMAClassName}{% if non-KOMA class
  \IfFileExists{parskip.sty}{%
    \usepackage{parskip}
  }{% else
    \setlength{\parindent}{0pt}
    \setlength{\parskip}{6pt plus 2pt minus 1pt}}
}{% if KOMA class
  \KOMAoptions{parskip=half}}
\makeatother
\usepackage{longtable,booktabs,array}
\newcounter{none} % for unnumbered tables
\usepackage{calc} % for calculating minipage widths
% Correct order of tables after \paragraph or \subparagraph
\usepackage{etoolbox}
\makeatletter
\patchcmd\longtable{\par}{\if@noskipsec\mbox{}\fi\par}{}{}
\makeatother
% Allow footnotes in longtable head/foot
\IfFileExists{footnotehyper.sty}{\usepackage{footnotehyper}}{\usepackage{footnote}}
\makesavenoteenv{longtable}
\setlength{\emergencystretch}{3em} % prevent overfull lines
\providecommand{\tightlist}{%
  \setlength{\itemsep}{0pt}\setlength{\parskip}{0pt}}
\usepackage{amsmath,amssymb,mathtools,needspace}
\xeCJKDeclareCharClass{CJK}{"2032,"0400->"04FF,"2160->"216B,"2460->"2473,"25A0,"25C7,"25CB,"25CF}
\IfFontExistsTF{Arial Unicode MS}{\xeCJKsetup{AutoFallBack=true}\setCJKfallbackfamilyfont{\CJKrmdefault}{Arial Unicode MS}}{}
\setcounter{secnumdepth}{0}
\setlength{\emergencystretch}{2em}
\allowdisplaybreaks
\usepackage{bookmark}
\IfFileExists{xurl.sty}{\usepackage{xurl}}{} % add URL line breaks if available
\urlstyle{same}
\hypersetup{
  pdftitle={幂法},
  colorlinks=true,
  linkcolor={Maroon},
  filecolor={Maroon},
  citecolor={Blue},
  urlcolor={Blue},
  pdfcreator={LaTeX via pandoc}}

\title{幂法}
\author{}
\date{}

\begin{document}
\maketitle

\subsubsection{9.2.1 幂法}\label{ux5e42ux6cd5}

在一些工程、物理问题中，通常只需要求出矩阵的按模最大的特征值（称为 \(A\) 的主特征值）和相应的特征向量，对于解这种特征值问题，应用幂法是合适的。

幂法是一种计算实矩阵 \(A\) 的主特征值的一种迭代法，它最大的优点是方法简单，对于稀疏矩阵较合适，但有时收敛速度很慢。

设实矩阵 \(A=(a_{ij})_n\) 有一个完全的特征向量组，其特征值为 \(\lambda_1,\lambda_2,\cdots,\lambda_n\)，相应的特征向量为 \(x_1,x_2,\cdots,x_n\)。已知 \(A\) 的主特征值是实根，且满足条件

\[
|\lambda_1|>|\lambda_2|\geq|\lambda_3|\geq\cdots\geq|\lambda_n|.
\tag{9.2.1}
\]

幂法的基本思想是任取一个非零的初始向量 \(v_0\)，由矩阵 \(A\) 构造一向量序列

\[
\left\{
\begin{aligned}
v_1&=Av_0,\\
v_2&=Av_1=A^2v_0,\\
&\quad\vdots\\
v_{k+1}&=Av_k=A^{k+1}v_0,\\
&\quad\vdots
\end{aligned}
\right.
\tag{9.2.2}
\]

称为迭代向量。由假设，\(v_0\) 可表示为

\[
v_0=a_1x_1+a_2x_2+\cdots+a_nx_n\quad\text{（设 }a_1\ne0\text{）},
\tag{9.2.3}
\]

于是

\[
\begin{aligned}
v_k&=Av_{k-1}=A^kv_0=a_1\lambda_1^kx_1+a_2\lambda_2^kx_2+\cdots+a_n\lambda_n^kx_n\\
&=\lambda_1^k\left[a_1x_1+\sum_{i=2}^na_1(\lambda_i/\lambda_1)^kx_i\right]=\lambda_1^k(a_1x_1+\varepsilon_k),
\end{aligned}
\]

其中 \(\varepsilon_k=\displaystyle\sum_{i=2}^na_1(\lambda_i/\lambda_1)^kx_i\)。由假设 \(|\lambda_i/\lambda_1|<1\ (i=2,3,\cdots,n)\)，故 \(\varepsilon_k\to0\ (k\to\infty)\)，从而

\[
\lim_{k\to\infty}\frac{v_k}{\lambda_1^k}=a_1x_1.
\tag{9.2.4}
\]

这说明序列 \(\dfrac{v_k}{\lambda_1^k}\) 越来越接近 \(A\) 的对应于 \(\lambda_1\) 的特征向量，或者说当 \(k\) 充分大时

\[
v_k\approx a_1\lambda_1^kx_1,
\tag{9.2.5}
\]

即迭代向量 \(v_k\) 为 \(\lambda_1\) 的特征向量的近似向量（除一个因子外）。

下面再考虑主特征值 \(\lambda_1\) 的计算。用 \((v_k)_i\) 表示 \(v_k\) 的第 \(i\) 个分量，则

\[
\frac{(v_{k+1})_i}{(v_k)_i}=\lambda_1\left\{\frac{a_1(x_1)_i+(\varepsilon_{k+1})_i}{a_1(x_1)_i+(\varepsilon_k)_i}\right\},
\tag{9.2.6}
\]

\begin{quote}
\textbf{校注（agent 补充，ch09a-E003）} 本页 \(v_k\) 第二行求和项及 \(\varepsilon_k\) 定义中的系数下标，原图均印为 \(a_1\)，已照录；与前一行 \(a_i\lambda_i^kx_i\) 的展开相比较，两处应为 \(a_i\) 才能保持恒等，属原书疑误。见\href{../assets/ch09a/review-p235-power-expansion.png}{放大原式}。
\end{quote}

\href{../../source-images/pdf-236.jpeg}{原页图像}

故

\[
\lim_{k\to\infty}\frac{(v_{k+1})_i}{(v_k)_i}=\lambda_1,
\tag{9.2.7}
\]

也就是说，两相邻迭代向量分量的比值收敛到主特征值。

这种由已知非零向量 \(v_0\) 及矩阵 \(A\) 的乘幂 \(A^k\) 构造向量序列 \(\{v_k\}\) 以计算 \(A\) 的主特征值 \(\lambda_1\)（利用式（9.2.7））及相应特征向量（利用式（9.2.5））的方法称为幂法。

由式（9.2.6）知，\((v_{k+1})_i/(v_k)_i\to\lambda_1\) 的收敛速度由比值 \(r=\lambda_2/\lambda_1\) 来确定，\(r\) 越小收敛越快，但当 \(r=\lambda_2/\lambda_1\approx1\) 时收敛可能就很慢。

总结上述讨论，有如下定理。

\textbf{定理 9.5} 设 \(A\in\mathbf R^{n\times n}\) 有 \(n\) 个线性无关的特征向量，主特征值 \(\lambda_1\) 满足

\[
|\lambda_1|>|\lambda_2|\geq|\lambda_3|\geq\cdots\geq|\lambda_n|,
\]

则对于任何非零初始向量 \(v_0\ (a_1\ne0)\)，式（9.2.4）、（9.2.7）成立。

设 \(A\) 的主特征值为实重根，即 \(\lambda_1=\lambda_2=\cdots=\lambda_r\)，且 \(|\lambda_r|>|\lambda_{r+1}|\geq\cdots\geq|\lambda_n|\)，又设 \(A\) 有 \(n\) 个线性无关的特征向量，\(\lambda_1\) 对应的 \(r\) 个线性无关特征向量为 \(x_1,x_2,\cdots,x_r\)，则由式（9.2.2），有

\[
v_k=A^kv_0=\lambda_1^k\left\{\sum_{i=1}^ra_ix_i+\sum_{i=r+1}^na_i\left(\frac{\lambda_i}{\lambda_1}\right)^kx_i\right\},\quad
\lim_{k\to\infty}\frac{v_k}{\lambda_1^k}=\sum_{i=1}^ra_ix_i\quad\text{（设 }\sum_{i=1}^ra_ix_i\ne0\text{）}.
\]

这说明当 \(A\) 的主特征值是实的重根时，定理 9.5 的结论还是正确的。

应用幂法计算 \(A\) 的主特征值 \(\lambda_1\) 及对应的特征向量时，如果 \(|\lambda_1|>1\)（或 \(|\lambda_1|<1\)），迭代向量 \(v_k\) 的各个不等于零的分量将随 \(k\to\infty\) 而趋于无穷（或趋于零），这样在用计算机计算时就可能``溢出''。为了克服这个缺点，就需要将迭代向量加以规范化。

设有一向量 \(v\ne0\)，将其规范化得到向量 \(u=\dfrac{v}{\max(v)}\)，其中 \(\max(v)\) 表示向量 \(v\) 的绝对值最大的分量。

在定理 9.5 的条件下幂法可这样进行：任取一初始向量 \(v_0\ne0\ (a_1\ne0)\)，构造向量序列

\[
\left\{
\begin{aligned}
v_1&=Au_0=Av_0,&u_1&=\frac{v_1}{\max(v_1)}=\frac{Av_0}{\max(Av_0)},\\
v_2&=Au_1=\frac{A^2v_0}{\max(Av_0)},&u_2&=\frac{v_2}{\max(v_2)}=\frac{A^2v_0}{\max(A^2v_0)},\\
&\quad\vdots&&\quad\vdots\\
v_k&=\frac{A^kv_0}{\max(A^{k-1}v_0)},&u_k&=\frac{A^kv_0}{\max(A^kv_0)}.
\end{aligned}
\right.
\]

由式（9.2.3），有

\[
Av_0=\sum_{i=1}^na_i\lambda_i^kx_i=\lambda_i^k\left[a_1x_1+\sum_{i=2}^na_i\left(\frac{\lambda_i}{\lambda_1}\right)^kx_i\right],
\tag{9.2.8}
\]

\begin{quote}
\textbf{校注（agent 补充，ch09a-E004）} 式（9.2.8）原图左端为 \(Av_0\)，右端方括号前为 \(\lambda_i^k\)，均照录。由 \(v_0=\sum_i a_ix_i\) 和 \(Ax_i=\lambda_ix_i\)，中间求和等于 \(A^kv_0\)，而提取的公共因子应为 \(\lambda_1^k\)；故有漏印幂次及因子下标疑误。见\href{../assets/ch09a/review-p236-9-2-8.png}{放大原式}。
\end{quote}

\href{../../source-images/pdf-237.jpeg}{原页图像}

\[
\begin{aligned}
u_k&=\frac{A^kv_0}{\max(A^kv_0)}
=\frac{\lambda_i^k\left[a_1x_1+\displaystyle\sum_{i=2}^na_i\left(\frac{\lambda_i}{\lambda_1}\right)^kx_i\right]}{\max\left[\lambda_1^k\left(a_1x_1+\displaystyle\sum_{i=2}^na_i\left(\frac{\lambda_i}{\lambda_1}\right)^kx_i\right)\right]}\\
&=\frac{\left[a_1x_1+\displaystyle\sum_{i=2}^na_i\left(\frac{\lambda_i}{\lambda_1}\right)^kx_i\right]}{\max\left[a_1x_1+\displaystyle\sum_{i=2}^na_i\left(\frac{\lambda_i}{\lambda_1}\right)^kx_i\right]}
\to\frac{x_1}{\max(x_1)}\qquad(k\to\infty).
\end{aligned}
\]

这说明规范化向量序列收敛到主特征值对应的特征向量。

同理，可得到

\[
v_k=\frac{\lambda_1^k\left[a_1x_1+\displaystyle\sum_{i=2}^na_i\left(\frac{\lambda_i}{\lambda_1}\right)^kx_i\right]}{\max\left[\lambda_1^{k-1}a_1x_1+\displaystyle\sum_{i=2}^na_i\left(\frac{\lambda_i}{\lambda_1}\right)^{k-1}x_i\right]},
\]

\[
\max(v_k)=\frac{\lambda_1\max\left[a_1x_1+\displaystyle\sum_{i=2}^na_i\left(\frac{\lambda_i}{\lambda_1}\right)^kx_i\right]}{\max\left[a_1x_1+\displaystyle\sum_{i=2}^na_i\left(\frac{\lambda_i}{\lambda_1}\right)^{k-1}x_i\right]}\to\lambda_1\qquad(k\to\infty).
\]

收敛速度由比值 \(r=\lambda_2/\lambda_1\) 确定。

总结上述讨论，有如下定理。

\textbf{定理 9.6} 设 \(A\in\mathbf R^{n\times n}\) 有 \(n\) 个线性无关的特征向量，主特征值 \(\lambda_1\) 满足 \(|\lambda_1|>|\lambda_2|\geq|\lambda_3|\geq\cdots\geq|\lambda_n|\)，则对于任意非零初始向量 \(v_0=u_0\ (a_1\ne0)\)，按下述方法构造的向量序列

\[
\left\{
\begin{aligned}
v_0&=u_0\ne0,\\
v_k&=Au_{k-1},\\
u_k&=\frac{v_k}{\max(v_k)}
\end{aligned}
\right.\qquad(k=1,2,\cdots),
\tag{9.2.9}
\]

有

\[
\lim_{k\to\infty}u_k=\frac{x_1}{\max(x_1)},\quad\lim_{k\to\infty}\max(v_k)=\lambda_1.
\]

\textbf{例 9.1} 用幂法计算 \(A=\begin{pmatrix}1.0&1.0&0.5\\1.0&1.0&0.25\\0.5&0.25&2.0\end{pmatrix}\) 的主特征值和相应的特征向量。

\textbf{解} 计算过程如表 9.1 所示。

下述结果是用 8 位浮点数字进行运算得到的，\(u_k\) 的分量值是舍入值。于是得到

\[
\lambda_1\approx2.536\,532\,3
\]

及相应特征向量 \((0.748\,2,\ 0.649\,7,\ 1)^{\mathrm T}\)。\(\lambda_1\) 和相应的特征向量真值（8 位数字）为

\[
\lambda_1=2.536\,525\,8,\quad\widetilde{x}_1=(0.748\,221\,16,\ 0.649\,661\,16,\ 1)^{\mathrm T}.
\]

\begin{quote}
\textbf{校注（agent 补充，ch09a-E005）} 页首 \(u_k\) 第一行分子方括号前原印 \(\lambda_i^k\)，分母对应因子为 \(\lambda_1^k\)，已照录；与 \(A^kv_0\) 的展开及下一行相消过程对照，分子因子应为 \(\lambda_1^k\)。见\href{../assets/ch09a/review-p237-u-k-numerator.png}{放大原式}。

\textbf{校注（agent 补充，ch09a-E006）} 本页 \(v_k\) 式分母中，原图 \(\lambda_1^{k-1}\) 仅直接乘 \(a_1x_1\)，随后为加号及求和项，已照录。由上一页 \(v_k=A^kv_0/\max(A^{k-1}v_0)\)，分母应把 \(\lambda_1^{k-1}\) 乘在整个 \(a_1x_1+\sum_{i=2}^na_i(\lambda_i/\lambda_1)^{k-1}x_i\) 上，故疑漏一层括号。见\href{../assets/ch09a/review-p237-v-k-denominator.png}{放大原式}。
\end{quote}

\href{../../source-images/pdf-238.jpeg}{原页图像}

\textbf{表 9.1}

{\def\LTcaptype{none} % do not increment counter
\begin{longtable}[]{@{}rcr@{}}
\toprule\noalign{}
\(k\) & \(u_k^{\mathrm T}\)（规范化向量） & \(\max(v_k)\) \\
\midrule\noalign{}
\endhead
\bottomrule\noalign{}
\endlastfoot
0 & \((1,1,1)\) & \\
1 & \((0.909\,1,\ 0.818\,2,\ 1)\) & \(2.750\,000\,0\) \\
5 & \((0.765\,1,\ 0.667\,4,\ 1)\) & \(2.558\,791\,8\) \\
10 & \((0.749\,4,\ 0.650\,8,\ 1)\) & \(2.538\,002\,9\) \\
15 & \((0.748\,3,\ 0.649\,7,\ 1)\) & \(2.536\,625\,6\) \\
16 & \((0.748\,3,\ 0.649\,7,\ 1)\) & \(2.536\,584\,0\) \\
17 & \((0.748\,2,\ 0.649\,7,\ 1)\) & \(2.536\,559\,8\) \\
18 & \((0.748\,2,\ 0.649\,7,\ 1)\) & \(2.536\,545\,6\) \\
19 & \((0.748\,2,\ 0.649\,7,\ 1)\) & \(2.536\,537\,4\) \\
20 & \((0.748\,2,\ 0.649\,7,\ 1)\) & \(2.536\,532\,3\) \\
\end{longtable}
}

\begin{center}\rule{0.5\linewidth}{0.5pt}\end{center}

\subsection{相关原页校注（agent 补充）}\label{ux76f8ux5173ux539fux9875ux6821ux6ce8agent-ux8865ux5145}

以下校注来自本节覆盖的原页，可能也涉及同页相邻小节；原文照录与校正说明分开保留。

\subsubsection{原书 PDF 页 238 的校注}\label{ux539fux4e66-pdf-ux9875-238-ux7684ux6821ux6ce8}

\href{../pages/pdf-238.md}{完整原页转录} · \href{../../source-images/pdf-238.jpeg}{原页图像}

校注记录：ch09a-E007。

\begin{quote}
\textbf{校注（agent 补充，ch09a-E007）} 本页``原点平移法''首段原印 \(r=\lambda_1/\lambda_2\)，已照录；前页及本页例 9.2 均用 \(r=\lambda_2/\lambda_1\)，与幂法误差项 \((\lambda_i/\lambda_1)^k\) 对照，此处疑将分子、分母颠倒。见\href{../assets/ch09a/review-p238-ratio.png}{放大原式}。
\end{quote}

\subsection{本节覆盖原页的校注记录}\label{ux672cux8282ux8986ux76d6ux539fux9875ux7684ux6821ux6ce8ux8bb0ux5f55}

下列记录按原页关联，含同页相邻小节的校注；计算时同时读取上文校注和完整原页。

\begin{itemize}
\tightlist
\item
  \texttt{ch09a-E003}：原书 PDF 页 235
\item
  \texttt{ch09a-E004}：原书 PDF 页 236
\item
  \texttt{ch09a-E005}：原书 PDF 页 237
\item
  \texttt{ch09a-E006}：原书 PDF 页 237
\item
  \texttt{ch09a-E007}：原书 PDF 页 238
\end{itemize}

\end{document}
